\documentclass[leqno]{jsarticle}
\usepackage{amssymb}
\usepackage{amsthm}
\usepackage{amsmath}
\usepackage{comment}
	%可換図式用
		%\usepackage[all]{xy}
	%重くなるので使わいときはコメント化しておく。
%定理環境 thmはあまり使わずpropをなるべく使う。
%主張は基本的に証明の中で使い、そのため番号を付けない。
%注意環境を付けてみた。
\newtheorem{thm}{定理}
\newtheorem{prop}[thm]{命題}
\newtheorem{cor}[thm]{系}
\newtheorem{lem}[thm]{補題}
\newtheorem{df}[thm]{定義}
\newtheorem{exam}[thm]{例}
\newtheorem{fact}[thm]{事実}
\newtheorem*{claim}{主張}
\newtheorem*{cau}{注意}
\renewcommand{\proofname}{{\bf 証明}}
%矢印たち。
%\toは\rightarrowと同じ。\getsは\leftarrowと同じ。同値は\iff
%上向き矢はuparrow逆はdownarrow
\newcommand{\To}{\Longrightarrow}
\newcommand{\Gets}{\Longleftarrow}
%黒板太字
\newcommand{\nn}{\mathbb{N}}
\newcommand{\zz}{\mathbb{Z}}
\newcommand{\qq}{\mathbb{Q}}
\newcommand{\rr}{\mathbb{R}}
\newcommand{\cc}{\mathbb{C}}
%無限大
\newcommand{\mgn}{\infty}
%ギリシャン
\newcommand{\dpst}{\displaystyle}
\newcommand{\ep}{\varepsilon}
\newcommand{\al}{\alpha}
\newcommand{\be}{\beta}
\newcommand{\de}{\delta}
\newcommand{\la}{\lambda}
\newcommand{\La}{\Lambda}
\newcommand{\ga}{\gamma}
\newcommand{\lil}{\la\in \La}
%商集合
\newcommand{\qt}[2]{#1/\! #2}
%enumerate環境
\renewcommand{\labelenumi}{{\rm (\arabic{enumi})}}
%以降alignじゃなくてenumerate を使え。
%なんか演算
\newcommand{\card}{\text{{\rm card}}}
\newcommand{\supp}{\text{{\rm supp}}}
%なんか演算２
\newcommand{\bd}{\text{{\rm Bdry}}}
\newcommand{\cl}{\text{{\rm CL}}}
\newcommand{\nai}{\text{{\rm Int}}}
%差集合は\setminusで統一する。等号の否定は\neq
%真部分集合は\subsetneqq　偏微分は\partial
%ディスプレイスタイル。\dfracでディスプレイ分数
%空集合は\Oを使う！！！！！！！！！！！！

\title{LOTSの位相の基本}
\author{電波通信}
\date{\today}
			\begin{document}
\maketitle
うろ覚えで書いたので標準的な用語と異なるかもしれません。
\section{定義など}
\begin{df}[線形順序]
集合$X$上の二項関係$R$が線形順序であるとは
\begin{enumerate}
\item $\forall a\in X\ aRa$
\item $\forall a, b\in X\ aRb\land bRa\To a=b$
\item $\forall a,b,c\in X\ aRb\land bRc\To aRc$
\item $\forall a,b\in X\ aRb\lor bRa$
\end{enumerate}
を充たすことを言う。慣例に従い線形順序を$\le$などと書き表すことにする。また、
\[
a<b\iff a\le b \land a\neq b
\]
と定義する。
\end{df}

\begin{df}[部分線形順序集合]
線形順序集合$(X,\le)$の部分集合$S$について$\le$を$S$に制限した関係を$\le $の$S$のへの制限といい、$\le_S$と書いてみたりするが、しばしば同じ記号$\le$を流用してしまう。大体文脈で分かるので問題ないのだが、区別する必要がある場合もある。このとき$(S,\le_S)$を$(X.\le)$の部分線形順序集合という。
\end{df}

\begin{df}[順序稠密性、順序の意味で稠密性]
線形順序集合$(X,\le )$の部分集合$A$が（順序の意味で）稠密であるとは$X$の任意の二元$a,b\in X$について$A$の元$c$が存在して
\[
a<c<b
\]
を充たすことを言う。これは順序位相における稠密性と同じことである。
\end{df}
\begin{df}
線形順序集合$(X,\le )$の元$a$が$S$の上界であるとは
\[
\forall x\in S\ x\le a
\]
を充たすときに言う。\par
線形順序集合$(X,\le )$の元$b$が$S$の下界であるとは
\[
\forall x\in S\ b \le x
\]
を充たすときに言う。
\end{df}

\begin{df}[$\sup, \inf, \max, \min$]
線形順序集合$(X,\le )$の元$s$が$S$の上界でありかつ
\[
\forall y<s\ \exists x\in S\ y<x\le s
\]
を充たすとき$s$を$S$の上限といい$\sup S$などと書き表す。特に$\sup S\in S$のとき$\sup S$は$S$の最大元といい$\max S$などと表す。
双対的に$\inf S$が定義される。\par
$(X,\le)$の元$t$が$S$の下限であるとは$t$は$S$の下界で
\[
\forall y>t\ \exists x\in S\ t\le x<y
\]
が成り立つときに言い、$\inf S$などと書く。$\inf S\in S$のとき$\min S$だとか書く。\par
上限と下限は常に存在するわけではないことに注意せよ。
\end{df}

\begin{cau}
$\sup S$は$S$の上界の最小値であり、$\inf S$は$S$の下界の最大限である。
\end{cau}

\begin{df}[順序完備性]
線形順序空間$(X,\le)$が順序完備であるとは、$X$の任意の非空な部分集合$S$について$\inf S$と$\sup S$が存在するときに言う。
\end{df}

\begin{df}[区間]
線形順序集合$(X,\le)$と$a,b\in X$について
\begin{align}
(a,b)=\{x\in X\mid a<x<b\}\\
(a,b]=\{x\in X\mid a<x\le b\}\\
[a,b)=\{x\in X\mid a\le x<b\}\\
[a,b]=\{x\in X\mid a\le x\le b\}
\end{align}
と定義する。それぞれ開区間、半開区間、閉区間という。また
\begin{align}
(a,\to)=\{x\in X\mid a<x \}\\
(\gets,a)=\{x\in X\mid x<a\}\\
[a\to),=\{x\in X\mid a\le x\}\\
(\gets,a]=\{x\in X\mid x\le a \}
\end{align}
も無限開区間、無限閉区間などという。
\end{df}

\begin{cau}
$b\le a$のとき、$(a,b)=\O$に注意せよ。
\end{cau}

\begin{df}
線形順序集合$(X,\le)$について$(X,\le )$の開区間及び無限開区間全体の族は開基の公理を充たす。この開基から生成される位相を$(X,\le)$の順序位相といい、ここでは$\mathfrak{O}_{\le}$と書くことにしよう。何らかの線形順序集合の順序位相と同相になっている位相空間をLOTS(Linearly Ordered TopolTogical Space)（線形順序位相空間）という。
\end{df}

\section{性質}

\begin{prop}
LOTSはハウスドルフ空間である。
\end{prop}
\begin{proof}
線形順序位相空間$(X,\le)$の二点$a,b$を任意に与える。このとき$a<b$としても一般性を失わない。\par
Case1:$(a,b)\neq \O$のとき;$c\in (a,b)$をとると、$a\in (\gets, c),\ b\in(c,\to)$であるからわかる。\par
Case2:$(a,b)=\O$のとき;この状況のもと、$(a,\to)=[b,\to), (\gets, b)=(\gets, a]$であるからわかる。
\end{proof}
\begin{cau}\normalfont
実は一般にLOTSは単調正規という結構強めの分離公理を充たす。どのくらいの強さかというと、単調正規ならば遺伝的族正規であるくらいの強さである。詳しくは\\
 R.W.Heath, D.J.Lutzer and P.L.zenor, \textit{Monotonically Normal spaces}, Trans.Amer.Math.Soc.178(1973), p481-493
を参照のこと
\end{cau}

\begin{prop}
LOTS$(X,\le)$の閉集合$A$に下限が存在するならば、$\inf A\in A$。上限についても同様。
\end{prop}
\begin{proof}
下限の定義から、$\inf A$は$A$の触点になるのでわかる。
\end{proof}

\begin{thm}
$(X,\le)$が順序完備であることと、$(X,\mathfrak{O}_{\le})$がコンパクトであることは同値である。
\end{thm}
\begin{proof}\ 
[順序完備$\To$コンパクト]\par
$\mathcal{U}$を$X$の任意の開被覆としよう。$X$には最小値と最大値が存在するので$a=\min X,\ b=\max S$しよう。そして
\[
S=\{x\in X\mid \text{$[a,x]$は$\mathcal{U}$の有限部分族で被覆される}\}
\]
とおく。$a\in S$なので$S\neq \O$である。また、$x\in S\To [a,x]\subset S$にも注意しよう。$\ s=\sup S$と置く。まずは$s\in S$を示そう。$[a,s)\subset S$となることはすぐにわかる。$s\in U$となる$U\in \mathcal{U}$を一つとる。そして$A=[a,s]\setminus U$は閉集合であるから$t=\sup A\in A$で$s\not\in A$より$t<s$が分かる。そして$t\in S$なので$[a,t]$は有限個の$\mathcal{U}$の元$V_1,V_2,\dots V_n$で被覆されるので結局$[a,s]$は$\mathcal{U}$の有限個の元$V_1,V_2,\dots V_n,U$で被覆されるから$s\in S$となる。\par
今、$[a,s]$を被覆する$\mathcal{U}$の有限個の元を$U_1,U_2,\dots,U_m$とし、$B=[s,b]\setminus \bigcup_{i=1}^mU_i$と定義する。$B$は閉集合で$s\notin B$である。もし$B=\O$ならば$p=\inf B\in B$が存在するから$p$を含む$\mathcal{U}$の元$O$を選べば$U_1,U_2,\dots U_m, O$は$[a,p]$を被覆するので$p\in S$だが他方$s<p$となので$s=\sup S$に反する。よって$B=\O$つまり$s=\max X$なので$X$は$\mathcal{U}$の有限個の元で被覆される。\par
[コンパクト$\To$順序完備]
$S$を$X$の非空な部分集合とする。このとき$X$の閉集合族
\[
\{\cl(S)\}\cup\{[a,\to)\mid a\in S\}
\]
は$X$の有限交叉族であるので$X$のコンパクト性から
\[
\cl(S)\cap \bigcap_{a\in S} [a,\to)\neq \O
\]
となる。この集合から元$s$をとると、$s$は$\dpst \bigcap_{a\in S} [a,\to)$の元であることから$S$の上界になり、$\cl(S)$に属していることから任意の$y<s$について
\[
S\cap (y,\to)\neq\O
\]
だが、$s$は$S$の上界なのでこれは
\[
S\cap (y,s]\neq \O
\]
ということであり、これは$s=\sup S$ということと同値である。よって$\sup S$が存在する。$\inf S$が存在することも同様である。
\end{proof}

\begin{cor}
順序数$\al$について$[0,\al]=\al+1$はコンパクトである。
\end{cor}

\section{部分集合について}
$(X,\le)$の部分集合$S$には、$\le_S$から誘導される順序位相$\mathfrak{A}(S)$と、LOTS$(X,\le)$の部分空間としての位相$\mathcal{B}(S)$の二種類が考えられるが、一般にこの二つは\textbf{一致しない。}ただし、$\mathcal{B}$の位相は$\mathcal{A}$より細かいことはよくわかる。つまり$\mathcal{A}(S)\subset \mathcal{B}(S)$。これは、$\le_S$による開区間を$(a,b)_{S}$と書き、$X$における開区間を$(a,b)_{X}$と書いたときに、$a,b,\in S$について
\[
(a,b)_S=(a,b)_X\cap S
\]
が成り立つことからわかる。
これら二つが一致する簡単な十分条件を二つ紹介しよう。

\begin{prop}
$S$が$X$の稠密な部分集合のとき、$\mathfrak{A}(S)=\mathfrak{B}(S)$
\end{prop}
\begin{proof}
稠密性から任意の$a,b\in X$について
\[
(a,b)_X\cap S=\bigcup_{s,t\in (a,b)_X\cap S}(s,t)_S
\]
となるのでわかる。
\end{proof}

\begin{prop}
$S$が$X$のコンパクト集合であるとき、$\mathfrak{A}(S)=\mathfrak{B}(S)$
\end{prop}
\begin{proof}
$(S,\mathfrak{B}(S))$がコンパクトで$(S,\mathfrak{A}(S))$はハウスドルフでかつ$\mathfrak{A}(S)\subset \mathfrak{B}(S)$から$\mathfrak{A}(S)=\mathfrak{B}(S)$が分かる。
\end{proof}

\begin{prop}
$S$が$X$の区間であるとき、$\mathfrak{A}(S)=\mathfrak{B}(S)$
\end{prop}
\begin{proof}
わかる。
\end{proof}



















			\end{document}



\begin{comment}
[集合のやつは$\mid$を使う。]
[真なる包含関係]
\subsetneqq
[可換図式]
\[\xymatrix{
&   & (2) \ar@{=>}[rd]&  &  &  & (5) \ar@{=>}[rd]&\\ 
& (1)\ar@{=>}[ru] \ar@{=>}[rd]&  &(4)\ar@{=>}[r]&(7)& (1)\ar@{=>}[ru] \ar@{=>}[rd]& & (7)&\\ 
&   & (3) \ar@{=>}[ur]&  &  &  &(6) \ar@{=>}[ur]&
}\]

[\bigin \endを使わない方法]

{\prop[aa] aaa}
で
\begin{prop}[aa]
aaa
\end{prop}
と同じ\df \thm \proof でも同様

[直和]
$\amalg$

[同値関係でよく使う波線]
\sim

[引数付きの新命令]
\newcommand{\prn}[1]{\|#1\|_p}

[番号のリセット]

\setcounter{equation}{0}


[字体の変更]

{\rm hogehoge}と使う。

\rm ローマン
\it イタリック
\sl 斜体

[supp]

\newcommand{\supp}{\text{{\rm supp}}}

[中央揃えの改行]

	\begin{eqnarray*}
		hoge1&=&hogehoge
		hoge2&=&hogehofe

	\end{eqnarray*}

&&の部分で揃えられる。


[\tag]
\begin{equation}
f(x) = ax^2 + bx + c \tag{1.2.3}
\end{equation}



[場合分け]

	\begin{cases}
		hoge1 & \text{状況１}\\
		hoge2 & \text{状況２}\\
		・
		・
		・
	\end{cases}



\end{comment}





